\documentclass[11pt]{article}

\usepackage{amsmath,amssymb,amsthm}
\usepackage{geometry}
\usepackage{mathtools}
\usepackage{bm}
\usepackage{xcolor}

\geometry{margin=1in}

\title{Stochastic Control Optimization for Autocorrelated Returns}
\author{}
\date{}

\begin{document}
\maketitle

\section{Model Setup}

We consider a process of returns $R_t$ composed of a forecast signal $F_t$ and white noise $\epsilon_t$:
\begin{equation}
    R_t = F_t + \epsilon_t,
\end{equation}
where $E[F_t] = 0$ and $E[\epsilon_t] = 0$, $\text{cov}(\epsilon_t, \epsilon_{t-k}) = \sigma_\epsilon^2 \delta_k$.

The forecast $F_t$ is modeled as an AR(1) process:
\begin{equation}
    F_t = \nu F_{t-1} + \eta_t, \quad |\nu| < 1,
\end{equation}
where $\nu = \text{corr}(F_t, F_{t-1})$ represents the implicit decay in the forecast. 
The autocovariance function of $F_t$ is given by:
\begin{equation}
    \rho(k) = \text{cov}(F_t, F_{t-k}) = \rho(0) \nu^{|k|}.
\end{equation}

We define a generating function $\phi(\lambda)$ for the autocovariance:
\begin{equation}
    \phi(\lambda) \coloneqq \sum_{k=0}^\infty \lambda^k \rho(k) = \frac{\rho(0)}{1 - \lambda \nu}.
\end{equation}

\section{Optimal Strategy Formulation}

We search for an optimal position process $P_t$ in the class of processes defined by a scale parameter $A$
    and an exponential decay parameter $\lambda$:
\begin{equation}
\begin{aligned}
    P_t &= A G_t, \\
    G_t &= F_t + \lambda G_{t-1} = \sum_{k=0}^\infty \lambda^k F_{t-k}.
\end{aligned}
\end{equation}
where $G_t$ is the exponentially weighted sum of past forecasts.

Equivalently, 
\begin{equation}
P_t = A F_t + \lambda P_{t-1}
\end{equation}

Note: $P_t$ is an $AR(2)$ process: 
\begin{equation}
    P_t =  A \eta_t + (\lambda + \nu) P_{t-1} - \lambda \nu P_{t-2}.
\end{equation}

\section{Derived Quantities and Costs}

We compute the unconditional (stationary) expectations of several quantities over a long trading horizon of $T$ steps. By stationarity, the expected sum over $T$ steps is simply $T$ times the unconditional expectation of a single step.

\subsection{Ideal P\&L (IPNL)}
The expected ideal P\&L is:
\begin{equation}
    E[\text{IPNL}] = E\left[ \sum_{t=1}^T P_t R_t \right] = E\left[ \sum_{t=1}^T P_t F_t \right] = A T \sum_{k=0}^\infty \lambda^k E[F_t F_{t-k}] = A T \phi(\lambda) = A T \frac{\rho(0)}{1 - \lambda \nu}.
\end{equation}

\subsection{Self-Self Impact Cost (SSI)}
SSI refers to the quadratic trade cost. The trades are $P_t - P_{t-1} = A(G_t - G_{t-1}) = A(F_t + (\lambda-1)G_{t-1})$.
\begin{equation}
    E[\text{SSI}] = \mu \frac{\sigma}{\text{ADV}} E\left[ \sum_{t=1}^T (P_t - P_{t-1})^2 \right]
\end{equation}
For the AR(1) case, this simplifies to:
\begin{equation}
    E[\text{SSI}] = A^2 T \frac{2 \mu \sigma}{\text{ADV}} \rho(0) \frac{1-\nu}{(1+\lambda)(1-\lambda\nu)}.
\end{equation}

\subsection{Self-Self Risk Cost (SSR)}
SSR is the quadratic holding cost (variance cost):
\begin{equation}
    E[\text{SSR}] = \gamma \sigma^2 E\left[ \sum_{t=1}^T P_t^2 \right] = A^2 T \gamma \sigma^2 E[G_t^2].
\end{equation}
For AR(1) forecasts:
\begin{equation}
    E[\text{SSR}] = A^2 T \gamma \sigma^2 \frac{\rho(0)}{1-\lambda^2} \frac{1+\lambda \nu}{1 - \lambda \nu}.
\end{equation}

\section{Utility Maximization}

The objective is to maximize the expected utility:
\begin{equation}
    U(A, \lambda) = E[\text{IPNL}] - E[\text{SSI}] - E[\text{SSR}].
\end{equation}
Ignoring the factor $T \rho(0)$ and substituting $b(\lambda) = \frac{1}{1-\lambda\nu}$, we can write:
\begin{equation}
    \frac{U(A, \lambda)}{T \rho(0)} = A b(\lambda) - A^2 a(\lambda),
\end{equation}
where
\begin{equation}
    a(\lambda) = \frac{1}{1 - \lambda \nu} \left( \frac{2 \mu \sigma}{\text{ADV}} \frac{1-\nu}{1+\lambda} + \gamma \sigma^2 \frac{1+\lambda\nu}{1-\lambda^2} \right).
\end{equation}

\subsection{Solving for $A$}
For a fixed $\lambda$, the optimal $A_*$ is:
\begin{equation}
    A_*(\lambda) = \frac{b(\lambda)}{2 a(\lambda)}.
\end{equation}

\subsection{Solving for $\lambda$}
Substituting $A_*$ back into the utility:
\begin{equation}
    U^*(\lambda) = \frac{b^2(\lambda)}{4 a(\lambda)}.
\end{equation}
Maximizing $U^*(\lambda)$ is equivalent to minimizing $V(\lambda) \coloneqq \frac{4 a(\lambda)}{b^2(\lambda)}$:
\begin{equation}
    V(\lambda) = 4(1 - \lambda \nu) \left[ \frac{2 \mu \sigma}{\text{ADV}} \frac{1-\nu}{1+\lambda} + \gamma \sigma^2 \frac{1+\lambda\nu}{1-\lambda^2} \right].
\end{equation}
The optimal $\lambda_*$ is found by setting the derivative to zero, which recaps to the quadratic solution shown in the notes:
\begin{equation}
    \lambda_* = \text{argmin}_\lambda V(\lambda).
\end{equation}

\end{document}
